\documentclass[10pt,a4paper]{article}
\usepackage{amsmath,amssymb,amsthm,graphicx}
\usepackage{hyperref}
% Editorial numbering adapter for ahl-129 "Hypocoercivity with Schur complements" (English paper).
% The cedram class natively declares the article's theorem family on one shared within-section
% counter; the article wrapper declares none of it. This adapter re-declares the family in the
% English names the journal prints, on one shared counter named cdrthm: the retained preamble's
% author line \AtBeginDocument{\newtheorem{assumption}[cdrthm]{Assumption}} (retained-preamble
% line 13, delivered verbatim inside the retained preamble span) requires that counter to exist
% at \begin{document}. Equations are numbered within sections exactly as the class prints them.
% The adapter adds no mathematics.
\newtheorem{cdrthm}{Theorem}[section]
\newtheorem{theo}[cdrthm]{Theorem}
\newtheorem{prop}[cdrthm]{Proposition}
\newtheorem{cor}[cdrthm]{Corollary}
\newtheorem{lem}[cdrthm]{Lemma}
\newtheorem{defi}[cdrthm]{Definition}
\newtheorem{rema}[cdrthm]{Remark}
\numberwithin{equation}{section}
% Retained preamble for this source, sliced MECHANICALLY from its own main TeX.
% Source file : full-original-native.tex (lines 1-154 of the file)
% The block below is verbatim author text; the wrapper class, document metadata and
% \begin{document} are omitted because the assembler supplies them.
% lines taken: 1, 2, 3, 4, 5, 6, 7, 8, 9, 10, 11, ... (154 lines total)
%~Mouliné par MaN_auto v.0.25.0 2022-03-23 11:09:47

\usepackage{mathrsfs}
\usepackage{amscd}
\usepackage{amssymb}






\AtBeginDocument{\newtheorem{assumption}[cdrthm]{Assumption}}




\newcommand{\rme}{\mathrm{e}}
\newcommand{\ri}{\mathrm{i}}
\newcommand{\rmc}{\mathrm{c}}
\newcommand{\rms}{\mathrm{s}}

\newcommand{\Id}{\mathrm{Id}}
\newcommand{\Stokes}{\mathrm{Stokes}}
\newcommand{\Ran}{\mathrm{Ran}}
\newcommand{\divl}{\mathrm{div}}
\newcommand{\Lang}{\mathrm{Lang}}
\newcommand{\Boltz}{\mathrm{Boltz}}
\newcommand{\LSI}{\mathrm{LSI}}


\DeclareMathOperator{\Ent}{Ent}
\newcommand{\cL}{\mathcal{L}}
\newcommand{\cLs}{\mathcal{S}}
\newcommand{\cLa}{\mathcal{A}}
\newcommand{\cB}{\mathcal{B}}
\newcommand{\cX}{\mathcal{X}}
\newcommand{\cD}{\mathcal{D}}
\newcommand{\cR}{\mathcal{R}}
\newcommand{\cH}{\mathcal{H}}
\newcommand{\cLM}{\mathcal{M}}
\newcommand{\cLN}{\mathcal{N}}
\newcommand{\cLU}{\mathcal{U}}
\newcommand{\clV}{\mathcal{V}}

\newcommand{\R}{\mathbb{R}}
\newcommand{\bbZ}{\mathbb{Z}}
\newcommand{\bbT}{\mathbb{T}}
\newcommand{\bbN}{\mathbb{N}}



\newcommand{\bfun}{\mathbf{1}}
\newcommand{\scrC}{\mathscr{C}}
\newcommand{\scrT}{\mathscr{T}}

\newcommand{\frkS}{\mathfrak{S}}
\newcommand{\frka}{\mathfrak{a}}
\newcommand{\frkb}{\mathfrak{b}}

\newcommand{\Schur}{\mathfrak{S}_0}
\newcommand{\Schurb}{\mathfrak{S}_1}

\newcommand{\cLham}{\mathcal{L}_{\rm ham}}
\newcommand{\cLFD}{\mathcal{L}_{\rm FD}}

\newcommand{\eps}{\varepsilon}
\newcommand{\subplus}{\textnormal{\texttt{+}}}
\renewcommand{\leq}{\leqslant}
\renewcommand{\geq}{\geqslant}
\renewcommand{\le}{\leqslant}
\renewcommand{\ge}{\geqslant}

\newcommand{\cLNH}{\cL_{\rm NH}}
%%%%%%%%%%%%%%%%%%%%%%%%%%%%%%%%%%%%%%%%%%%%%%%%%%%%%%%%%%%
\graphicspath{{./figures/}}

\newcommand*{\mk}{\mkern -1mu}
\newcommand*{\Mk}{\mkern -2mu}
\newcommand*{\mK}{\mkern 1mu}
\newcommand*{\MK}{\mkern 2mu}

%\hypersetup{urlcolor=purple, linkcolor=blue, citecolor=red}

\newcommand*{\romanenumi}{\renewcommand*{\theenumi}{\roman{enumi}}}
\newcommand*{\Romanenumi}{\renewcommand*{\theenumi}{\Roman{enumi}}}
\newcommand*{\alphenumi}{\renewcommand*{\theenumi}{\alph{enumi}}}
\newcommand*{\Alphenumi}{\renewcommand*{\theenumi}{\Alph{enumi}}}
\let\oldtilde\tilde
\renewcommand*{\tilde}[1]{\mathchoice{\widetilde{#1}}{\widetilde{#1}}{\oldtilde{#1}}{\oldtilde{#1}}}
\let\oldforall\forall
\renewcommand*{\forall}{\mathrel{\oldforall}}
%%%%%%%%%%%%%%%%%%%%%%%%%%%%%%%%%%%%%%%%%%%%%%%%%%%%%%%%%%%%%%

\usepackage{cleveref}
\crefname{originalproposition}{original proposition}{original propositions}
\Crefname{originalproposition}{Original proposition}{Original propositions}
\makeatletter
\newlabel{prop:estimate_cLa^2_invBBstar}{{3.5}{source}{}{sourcepointer-1}{}}
\newlabel{prop:estimate_cLa^2_invBBstar@cref}{{[originalproposition][0][]3.5}{[1][0][]source}{}{}{}}
\makeatother
\crefname{originallemma}{original lemma}{original lemmas}
\Crefname{originallemma}{Original lemma}{Original lemmas}
\makeatletter
\newlabel{lem:VillaniA24}{{3.7}{source}{}{sourcepointer-2}{}}
\newlabel{lem:VillaniA24@cref}{{[originallemma][0][]3.7}{[1][0][]source}{}{}{}}
\makeatother
\title{Improved constants in the H2 estimate for the Langevin dynamics: C = 2 and logarithmic dimension scaling of C'}
\author{Étienne Bernard, Max Fathi, Antoine Levitt, Gabriel Stoltz}
\date{doi:10.5802/ahl.129; CC BY 4.0}
\begin{document}
\maketitle
\setcounter{section}{3}
\setcounter{cdrthm}{8}
\setcounter{equation}{0}
\noindent\textit{Editorial note.} The counted unit is the author's labelled proposition prop:scaling\_dimension\_Langevin\_improved at source lines 899-908 with the author's own proof environment immediately following at source lines 912-946; the pairing is the ordinary adjacent one (the prose note at source line 910 between them is not delivered).\par\medskip
\noindent\textit{Editorial note.} Journal numbering: the proposition prints as Proposition 3.9 in the publisher PDF, the ninth theorem-family item of Section 3 under cedram's shared family counter (Assumption 3.1, Assumption 3.2, Theorem 3.3, Corollary 3.4, Proposition 3.5, Lemma 3.6, Lemma 3.7, Definition 3.8 precede it in the source at lines 581, 593, 607, 721, 736, 791, 802 and 886). Erratum against the census: the census v3.11 printed field for this unit says 3.2; the publisher PDF anchor governs and the erratum is reported to the duty line. The card seeds the section counter with 3 and the shared family counter cdrthm with 8, so the counted proposition prints as 3.9 exactly as in the journal.\par\medskip
\noindent\textit{Editorial note.} Renumbering declarations: the three context blocks follow the counted proof in the excerpt and one of them is a numbered theorem environment -- the log-Sobolev definition prints as Definition 3.10 in the delivery baseline (journal 3.8) because the counted proposition occupies the ninth shared-counter slot ahead of it. The counted statement and proof carry three numbered equation environments in the delivery baseline: eq:condition\_Hessian\_LSI inside the statement and the sharper bound and the multline inside the proof print as (3.1), (3.2) and (3.3) respectively (journal (3.15) for the first; the two proof equations are unlabelled in the source and print among Section 3's journal numbering). The context equation eq:quadratic\_kinetic\_energy prints as (3.1) ahead of them and eq:Bochnereq\_CS as (3.2); the delivery baseline renumbering is declared here for the record.\par\medskip
\noindent\textit{Editorial note.} Citations: no delivered span (statement, proof, context) carries a \textbackslash{}cite, so the card ships no bibliography (bibliography\_paths empty, backend none). The two external references are carried as stubs: Proposition prop:estimate\_cLa\textasciicircum{}2\_invBBstar (journal Proposition 3.5, source lines 736-750), load-bearing as the counted conclusion's target and the delegated final assembly, stubbed because the sibling card S9-ahl129-estimate counts it; Lemma lem:VillaniA24 (journal Lemma 3.7, source lines 802-809), non-load-bearing, mentioned only to be set aside.\par\medskip
\noindent\textit{Editorial note.} Preamble composition: three ranges -- the editorial numbering adapter first (family re-declaration on the cdrthm anchor counter, English names, numberwithin), then the retained preamble's provenance comments (lines 1-6) and its body (lines 8-99). Line 7, the active cedram documentclass, is excluded because the wrapper supplies the class; lines 100-160 (title, alttitle, authors, keywords, abstracts and front matter) are excluded as publisher class-metadata commands that carry no mathematical content. hyperref is declared through extra\_packages for the retained hypersetup line; the wrapper's automatic amsmath, amssymb, amsthm and graphicx load precedes the preamble, and the retained amssymb usepackage line is an inert duplicate. All notation macros used by the delivered spans (rme, rmc, cD, cX, cB, R, eps, subplus, bbN, cLa, cLs, LSI, Ent, norm) are defined in the retained preamble itself; no mathematical macro, symbol, statement or proof step was added, changed or removed.\par\medskip
\noindent\textit{Editorial note.} Occupancy and claim evidence: claimkey\_index carries no ahl.129 key; doi\_public\_occupancy.py --doi 10.5802/ahl.129 returns PASS with occupied=0; the census v3.11 row for this unit reads pairing own with occupied false.\par\medskip
\noindent\textit{Allowed standard prerequisites (editorial).}\begin{itemize}
\item The standing Langevin setting of Section 3.1: D the position space with the potential V, beta the inverse temperature, nu the Gibbs measure e\textasciicircum{}\{-beta V\} dq / Z on D, and the Langevin generator with its symmetric and antisymmetric parts. the Section 3 opening prose around source lines 560-580, outside the counted spans. Fixes D, beta, V and nu in the counted statement and proof; the prose also fixes the notations ||nabla\_q u||\_\{L\textasciicircum{}2(nu)\} and nabla\_q\textasciicircum{}* used throughout.
\item The two matrix/vector norms used by the Hessian growth condition: for M in R\textasciicircum{}\{d x d\} the operator norm ||M||\_\{B(l\textasciicircum{}2)\} = sup\_\{|xi| <= 1\} |M xi|, and the l\textasciicircum{}inf norm |xi|\_inf = max(|xi\_1|, ..., |xi\_d|). source lines 893-897, outside the counted spans (the prose also cites Ledoux [Led99] for LSI background and is therefore not delivered as a context span). Defines the norms appearing in the numbered assumption equation eq:condition\_Hessian\_LSI of the counted statement.
\end{itemize}\medskip
\ifdefined\hypertarget\hypertarget{sourcepointer-1}{}\fi
\noindent\textit{Source reference.} Original proposition 3.5: sources/ahl-129-hypocoercivity-with-schur-complements-2022/full-original-native.tex, source lines 736--750. The counted statement concludes that the conclusion of this proposition holds true with explicit constants C = 2 and the logarithmic C-prime, and the counted proof closes by delegating to the end of that proposition proof (proof line 945). The proposition is a distinct unit of the same source, counted by the sibling card S9-ahl129-estimate (statement source lines 736-750, proof source lines 765-789), so this card stubs it with its journal number rather than duplicating its statement or proof.\par\medskip
\ifdefined\hypertarget\hypertarget{sourcepointer-2}{}\fi
\noindent\textit{Source reference.} Original lemma 3.7: sources/ahl-129-hypocoercivity-with-schur-complements-2022/full-original-native.tex, source lines 802--809. Mentioned in the counted proof only to be set aside (proof line 917: We cannot use Lemma lem:VillaniA24 to control the last term, but use instead the classical entropy inequality); the lemma is never invoked, so the stub carries its journal number 3.7 for reference resolution only. It is not counted by this card or by a sibling card.\par\medskip
\begin{prop}\label{prop:scaling_dimension_Langevin_improved}
Consider the quadratic kinetic energy~\eqref{eq:quadratic_kinetic_energy}. Assume that $\rme^{-\beta V} \in L^1(\cD)$, that $\nu$ satisfies a logarithmic Sobolev inequality, and that
\begin{equation}\label{eq:condition_Hessian_LSI}
\forall q \in \cD, \qquad \left\|\nabla^2V(q)\right\|_{\cB\left(\ell^2\right)} \leq c_3\left(1 +\left|\nabla V(q)\right|_{\infty}\right).
\end{equation}
Then the conclusion of Proposition~\ref{prop:estimate_cLa^2_invBBstar} holds true with
\[
C = 2, \qquad C' = 2\left(c_3 + \frac{1}{2c_3 C_{\LSI}}\left[\log d + \log\left(\max_{i=1,\,\dots,\,d} \int{\rme^{2c_3 C_{\LSI}\left|\partial_i V\right|} \, d\nu}\right)\right]\right).
\]
\end{prop}

\begin{proof}
The proof follows the same general approach as that of Proposition~\ref{prop:estimate_cLa^2_invBBstar}, but with the extra assumption on $V$, \eqref{eq:Bochnereq_CS} can be replaced by
\begin{equation}
\left\|\nabla_q^2 u\right\|_{L^2(\nu)}^2 \leq \left\|\nabla_q^{*}\nabla_q u \right\|_{L^2(\nu)}^2+ c_{3}\left\|\nabla_q u\right\|_{L^2(\nu)}^2 + c_{3}\int{\left|\nabla_q u\right|^2\left|\nabla V\right|_{\infty}d\nu}.
\end{equation}
We cannot use Lemma~\ref{lem:VillaniA24} to control the last term in the above equation, but use instead the following classical entropy inequality: for $g$ bounded measurable,
\[
\int_\cD f^2g \, d\nu \leq \Ent_{\nu}\left(f^2\right) + \|f\|_{L^2(\nu)}^2\log \left(\int_\cD \rme^{g} \, d\nu \right).
\]
This inequality is an immediate consequence of the fact that the entropy is the Legendre transform of the log-Laplace functional, with respect to the duality between functions and measures, i.e.
\[
\Ent_{\nu}\left(f^2\right) = \sup \left\{ \int_\cD f^2 g \, d\nu - \|f\|_{L^2(\nu)}^2\log \left(\int_\cD \rme^{g} \, d\nu\right), \; g \text{ bounded measurable} \right\}.
\]
As a consequence, we obtain a bound of the form
\[
\int_\cD \left|\nabla_q u\right|^2\left|\nabla V\right|_{\infty} \, d\nu \leq \frac1c \Ent_{\nu}\left(\left|\nabla_q u\right|^2\right) + \frac1c \left\|\nabla_q u\right\|_{L^2(\nu)}^2 \log \left(\int_\cD \rme^{c \left|\nabla V\right|_{\infty}} \, d\nu\right),
\]
with $c>0$ a constant to be adjusted later. When $\nu$ satisfies a logarithmic Sobolev inequality, the first factor on the right-hand side can be controlled by $C_{\LSI}\|\nabla^2_q u\|^2_{L^2(\nu)}$, using the inequality $\left|\nabla |\nabla f|\right|^2 \leq |\nabla^2 f|^2$. The second factor can be bounded as
\[
\log \left(\int_\cD \rme^{c|\nabla V|_{\infty}} \, d\nu\right) \leq \log \left(\sum_{i=1}^d \int_\cD \rme^{c\left|\partial_{q_i} V\right|_{\infty}} \, d\nu\right) \leq \log d + \log\left(\max_{i=1,\,\dots,\,d} \int_\cD \rme^{c\left|\partial_{q_i} V\right|} \, d\nu \right).
\]
We therefore obtain the inequality
\begin{multline}
\left\|\nabla^2 u\right\|_{L^2(\nu)}^2 \leq \left\|\nabla_q^{*}\nabla_q u \right\|_{L^2(\nu)}^2+ c_{3}\|\nabla_q u\|_{L^2(\nu)}^2 
+ \frac{c_{3}C_{\LSI}}{c} \left\|\nabla^2 u\right\|_{L^2(\nu)}^2 \\+ \frac1c \left\|\nabla_q u\right\|_{L^2(\nu)}^2 \left[\log d + \log\left(\max_{i=1,\,\dots,\,d} \int_\cD \rme^{c\left|\partial_i V\right|} \, d\nu \right)\right].
\end{multline}
Taking $c = \varepsilon^{-1}c_3 C_{\LSI}$ leads to
\begin{gather*}
\left\|\nabla^2 u\right\|_{L^2(\nu)}^2 \leq \frac{1}{1-\varepsilon}\left\|\nabla_q^{*}\nabla_q u \right\|_{L^2(\nu)}^2 + L'_{d,\, \varepsilon}\left\|\nabla_q u\right\|_{L^2(\nu)}^2,
\\
\intertext{with}
L'_{d,\,\varepsilon} := \frac{1}{1-\varepsilon}\left(c_3 + \frac{\varepsilon}{c_3 C_{\LSI}}\left[\log d + \log\left(\max_{i=1,\,\dots,\,d} \int_\cD \rme^{\varepsilon^{-1}c_3 C_{\LSI}\left|\partial_i V\right|} \, d\nu \right)\right]\right).
\end{gather*}
We next set $\eps = 1/2$ and follow the same reasoning as in the end of the proof of Proposition~\ref{prop:estimate_cLa^2_invBBstar} to conclude.
\end{proof}

\begin{equation}\label{eq:quadratic_kinetic_energy}
U(p) = \sum_{i=1}^d u(p_i), \qquad u(p) = \frac{p^2}{2m}.
\end{equation}

\begin{multline}\label{eq:Bochnereq_CS}
\left\|\nabla_q^2 u\right\|_{L^2(\nu)}^2 \\
\leq \left\|\nabla_q^{*}\nabla_q u \right\|_{L^2(\nu)}^2+ c_{3}\sqrt{d}\left\|\nabla_q u\right\|_{L^2(\nu)}^2 + c_{3}\left\|\nabla_q u\right\|_{L^2(\nu)}\big\| \left|\nabla_{q}u\right| \, \left| \nabla_{q}V\right| \big\|_{L^2(\nu)}.
\end{multline}

\begin{defi}
A probability measure~$\nu$ on $\cD$ is said to satisfy a logarithmic Sobolev inequality with constant $C_{\LSI}$ if
\begin{multline*}
\forall f \in C_\rmc^\infty(\cD), \\ \Ent_{\nu}\left(f^2\right) := \int_{\cD} f^2\log f^2 \, d\nu - \|f\|_{L^2(\nu)}^2\log \|f\|_{L^2(\nu)}^2 \leq C_{\LSI}\|\nabla f\|_{L^2(\nu)}^2.
\end{multline*}
\end{defi}

\end{document}
